\section{De GPT-2 a GPT-3: convertir la intuición en una Scaling Hypothesis verificable}

Mann considera que ImageNet en 2015 y la posterior aparición de GPT-2 fueron un field signal decisivo: las redes neuronales, que durante mucho tiempo se habían visto solo como una academic promise, empezaron a mostrar capacidades transferibles en tareas reales. GPT-2 sugería que «hacer next-token prediction sobre texto de internet a gran escala podría producir capacidades amplias», y GPT-3 usó experimentos a múltiples escalas para verificar si el crecimiento de los parámetros, los datos y el compute mejoraba de forma estable la loss y la few-shot performance.

\subsection{Observation, hypothesis, pilot y frontier run}

La \term{scaling law} (ley de escalamiento) es una regularidad empírica entre el desempeño o la loss de un modelo y el compute, el dataset size y el parameter count. Su valor de ingeniería no está en trazar a posteriori una línea elegante, sino en usar runs más pequeños antes del entrenamiento más grande para estimar el beneficio, comparar recipes, configurar la mezcla data/model y precisar cuándo la nueva evidencia indica que el regime ha cambiado.

\lecturefigure{02-scaling-hypothesis.png}{La scaling law convierte el field signal en una hypothesis, en pilot runs y en un frontier commitment predecible.}{Subtítulos locales 01:00--05:40; artículos de GPT-2/GPT-3 y de scaling laws.}

\begin{knowledgebox}{Lectura de la figura: antes del run más grande debe existir una predicción}
El field signal aporta la dirección; la hypothesis precisa «cómo cambia la loss al aumentar qué variables»; los pilot runs deben abarcar varias escalas y mantener una recipe comparable, y el frontier run, antes de que se apruebe el presupuesto, deja por escrito la expected curve, la confidence y la stop condition. Si el equipo solo tiene un run grande, sin holdout prediction, es difícil determinar si el éxito se debe al scaling, a la data cleanup, a un architecture change o a una interpretación selectiva.
\end{knowledgebox}

\teachervoice{Mann enfatiza que no siguió la trayectoria tradicional de un PhD y, aun así, participó en GPT-3 mediante data engineering, análisis y architecture experiments. Nota de clase: las contribuciones decisivas de la frontier research no provienen solo de proponer fórmulas nuevas, sino también de construir datos, experimentos y sistemas de medición confiables.}

\subsection{Por qué en ese momento mucha gente no lo creía}

El escepticismo no era absurdo: las tecnologías históricamente suelen seguir una sigmoide, los modelos existentes podían ser costosos y difíciles de desplegar, y la experiencia de la época de BERT/T5 privilegiaba el fine-tuning task-specific. El problema está en confundir «la economía bajo la recipe y el hardware actuales» con «el límite permanente de la capacidad». Una mejor forma de debatir es exigir predicciones entre escalas, residual analysis y límites de extrapolación, en lugar de rechazar toda la hipótesis con una demo costosa o con un único fracaso.

\begin{warningbox}{La scaling law no significa «basta con agregar compute para tener éxito»}
La architecture, el optimizer, la token mixture, la data contamination, la context length, la precision, el parallelism y el post-training pueden modificar la curva. Una power law solo describe la relación empírica dentro del intervalo observado; al extrapolar entre modalities, entre recipes o a lo largo de varios órdenes de magnitud, conviene ampliar la uncertainty y estar preparado para encontrar un nuevo bottleneck.
\end{warningbox}

\subsection{Resumen de la sección}

El carácter científico de la scaling hypothesis proviene de que es predecible y falsable. El field signal establece la dirección, los pilot runs a múltiples escalas aportan la curva y el frontier run pone a prueba la extrapolación; todo resultado debe contrastarse con la predicción preregistrada y con los residuals.
